\documentclass[11pt]{article}
\usepackage[utf8]{inputenc}
\usepackage[T1]{fontenc}
\usepackage{amsmath,amssymb,amsthm}
\usepackage{hyperref}
\usepackage{listings}

\title{Category Theory for Dan, Mike, and Ulrik}
\author{Groupoid Infinity}
\date{May 2026}

\newcommand{\catC}{\mathcal{C}}
\newcommand{\catD}{\mathcal{D}}
\newcommand{\op}{^\text{op}}
\newcommand{\homtype}[3]{\text{hom}(#1, #2, #3)}
\newcommand{\dan}[0]{\textbf{Dan}}
\newcommand{\mike}[0]{\textbf{Mike}}
\newcommand{\ulrik}[0]{\textbf{Ulrik}}
\newcommand{\stt}[0]{\textbf{STT}}
\newcommand{\dtt}[0]{\textbf{DTT}}

\begin{document}

\maketitle

\begin{abstract}
This article describes the design and implementation of the categorical libraries across the \dan{} presentation front-end and the \mike{} (DTT) and \ulrik{} (STT) typechecking environments. We focus on the transition from first-order Directed Type Theory (\mike{}), which utilizes ends, coends, and path-induction primitives, to the minimal primitives of Simplicial Type Theory (\ulrik{}), where all functions are functorial and natural transformations/Yoneda maps are formulated directly using dependent functions and reflexivity paths. We also outline the compilation and verification workflow via the \texttt{lift} utility.
\end{abstract}

\tableofcontents
\newpage

\section{Introduction}

The categorical libraries in our system are located in three corresponding directories:
\begin{itemize}
    \item \texttt{library/dan/}: Combinatorial definitions and simplicial complex shapes.
    \item \texttt{library/mike/}: Precategories, functors, natural transformations, adjunctions, and Yoneda coends formulated in Directed Type Theory (\dtt{}).
    \item \texttt{library/ulrik/}: Categories, functors, natural transformations, equivalences, and Yoneda maps formulated in Simplicial Type Theory (\stt{}).
\end{itemize}

\section{Definitional Duality}

In both \ulrik{} and \mike{} engines, the opposite category operation $(\cdot)\op$ is a definitional involution on categories:
\begin{equation}
    (\catC\op)\op \equiv \catC
\end{equation}
Furthermore, the dual of a natural transformation is definitionally involutive:
\begin{equation}
    (\alpha\op)\op \equiv \alpha
\end{equation}
This allows us to obtain dual category-theoretic statements "for free" upon evaluation in the typechecker, eliminating the need for explicit transport along isomorphisms.

\section{Directed Type Theory Library (Mike)}

The DTT engine (\mike{}) treats categories as first-order parameter structures and represents precategory structures using the $\homtype{\catC}{x}{y}$ primitive.
\begin{itemize}
    \item \textbf{Morphism Composition}: Defined via path-induction $J$:
    \begin{equation}
        g \circ f \equiv J(w_1.w_2.\homtype{\catC}{x}{w_2}, z, f, y, z, g)
    \end{equation}
    \item \textbf{Integrals}: DTT supports ends and coends natively. Natural transformations between functors $F$ and $G$ are defined as the end over the homs of the components:
    \begin{equation}
        \text{Nat}(F, G) \equiv \int_{w:\catC} \homtype{\catD}{F(w)}{G(w)}
    \end{equation}
    \item \textbf{Yoneda Coend}: The Yoneda lemma in DTT is expressed using coends and tensor products:
    \begin{equation}
        \text{YonedaCoend}(x, z) \equiv \int^{y:\catC} (\homtype{\catC}{x}{y} \otimes \homtype{\catD}{z}{F(y)})
    \end{equation}
    The covariant Yoneda map is defined using contravariant path induction:
    \begin{equation}
        \text{yoneda\_map} \equiv \lambda e. \text{let } \langle y, t \rangle := e \text{ in let } f \otimes u := t \text{ in } J_\text{contra}(w.F(w, z), u, y, f)
    \end{equation}
\end{itemize}

\section{Simplicial Type Theory Library (Ulrik)}

In Simplicial Type Theory (\stt{}) as implemented in \ulrik{}, we adopt a deliberately minimalistic set of primitives following the LICS 2025 Gratzer-Weinberger-Buchholtz paper. Primitives such as ends, coends, tensors, and path-induction operators ($J_\text{cov}$, $J_\text{contra}$) are omitted.

\begin{enumerate}
    \item \textbf{Functors}: All functions between types are automatically functorial in STT. Hence, a functor is simply represented as a function $F : A \to B$.
    \item \textbf{Natural Transformations}: A natural transformation between $F$ and $G$ is represented as a dependent function mapping objects to paths in $B$:
    \begin{equation}
        \text{Nat}(F, G) \equiv (w : A) \to \homtype{B}{F(w)}{G(w)}
    \end{equation}
    The identity natural transformation for a functor $F$ is defined using the constant reflexivity path $\lambda t. F(w)$ for each component:
    \begin{equation}
        \text{id\_nat} \equiv \lambda w. \lambda t. F(w)
    \end{equation}
    \item \textbf{Presheaves and Yoneda Map}: A presheaf $F$ on category $C$ is a function $C \to U$, where $U$ is the universe of spaces (types). Because morphisms in $U$ are functions ($\to$), natural transformations between presheaves $F$ and $G$ have type:
    \begin{equation}
        \text{PshNat}(F, G) \equiv (w : C) \to F(w) \to G(w)
    \end{equation}
    The representable functor for an object $c$ is $\text{Yo}(c) \equiv \lambda c'. \homtype{C}{c}{c'}$. The Yoneda map has type:
    \begin{equation}
        \text{yoneda\_map} : \text{PshNat}(\text{Yo}(c), F) \to F(c)
    \end{equation}
    This is defined by evaluating the natural transformation at the constant identity path:
    \begin{equation}
        \text{yoneda\_map}(\eta) \equiv \eta(c, \lambda t. c)
    \end{equation}
\end{enumerate}

\section{The Lifting and Verification Workflow}

The operational layer (\dan{}) allows for a presentation-oriented description of simplicial structures. Using the `lift` binary, these definitions are automatically translated into type-theoretic definitions:

\begin{lstlisting}[language=bash]
$ lift -stt library/dan/simplex.dan
$ lift -dirtt library/dan/simplex.dan
\end{lstlisting}

Verification is performed independently using the two separate typecheckers:
\begin{lstlisting}[language=bash]
$ dune exec ulrik library/ulrik/book.ulrik
$ dune exec mike library/mike/book.mike
$ dune exec dan library/dan/book.dan
\end{lstlisting}

\begin{thebibliography}{9}
\bibitem{RS2017}
Emily Riehl and Michael Shulman.
\textit{A type theory for synthetic \(\infty\)-categories}.
Higher Structures, 2017.
\bibitem{GWB2025}
Daniel Gratzer, Jonathan Weinberger, and Ulrik Buchholtz.
\textit{The Yoneda embedding in simplicial type theory}.
LICS 2025.
\end{thebibliography}

\end{document}
